\documentclass[10pt,conference]{IEEEtran}
\usepackage{cite}
\usepackage{amsmath,amssymb}
\usepackage[hidelinks]{hyperref}
\usepackage{booktabs}
\usepackage{array}
\input{claim_macros}

\newcommand{\Allow}{\mathsf{ALLOW}}
\newcommand{\Deny}{\mathsf{DENY}}
\newcommand{\App}{\mathsf{app}}
\newcommand{\Rej}{\mathsf{rej}}
\newcommand{\Unsup}{\mathsf{unsup}}
\newcommand{\Fail}{\mathsf{fail}}
\newcommand{\Core}{\mathsf{Core}}
\newcommand{\Eval}{\mathcal{E}}

\title{Supplementary Material for\\LatticeGuard: Law-Guided Regression Oracles for Access-Control Evaluators}
\author{Anonymous}

\begin{document}
\maketitle

\section{Evidence Boundary}
The main paper carries the scientific claim: access-control laws can be compiled into applicability-checked metamorphic obligations, and those obligations give a safer regression oracle than point examples, differential disagreement, or unrestricted mutation.  This supplement expands the evidence boundary behind that claim.  It should be read as a companion to the method, not as a place where the contribution is hidden.  The central definitions, results, limitations, and novelty argument remain in the main paper.

The evidence boundary is deliberately conservative.  LatticeGuard is not a bug-counting study and does not claim that Casbin, Cedar, or OPA are bug-free.  A clean primary run is meaningful only after three other facts are visible: invalid transformations are rejected before accounting, unsupported language fragments remain outside the denominator, and the counterexample path is exercised by semantic-drift controls.  This supplement therefore focuses on the questions that a skeptical reviewer is likely to ask after seeing a zero-primary-failure result.

\begin{table}[t]
\centering
\caption{How to read each evidence class.}
\label{tab:reading}
\begin{tabular}{p{0.24\linewidth}p{0.28\linewidth}p{0.28\linewidth}}
\toprule
Class & Supports & Does not imply \\
\midrule
Applicable obligations & Declared laws hold for evaluated fragments & Complete language coverage \\
Rejected candidates & Invalid oracle rows remain visible & Backend failures \\
Unsupported fragments & Adapter boundaries are explicit & Evidence about unsupported semantics \\
Seeded drifts & Failure and minimization paths are live & Production defects \\
Bounded law checks & Relation predicates survive finite falsification & Universal language proof \\
Trace replay & Reported counts can be recomputed & Scale as the contribution \\
\bottomrule
\end{tabular}
\end{table}

This reading rule avoids a common trap in regression-oracle papers.  More rows are not automatically stronger evidence.  Rows become evidence only after a law predicate says why a before/follow-up pair should preserve or order a decision.  Otherwise a generator can create thousands of syntactically valid edits whose expected relation is false.  The resulting pass rate may look precise while measuring the generator's accidental bias.  LatticeGuard makes the admission step part of the result.

\section{Formal Accounting Semantics}
Let $c$ be a generated candidate for law $\ell$.  The candidate contains the original policy slice, the follow-up slice, the original request, the follow-up request, and the declared relation family.  Three predicates decide whether it can enter the primary denominator: a rejection predicate $\rho_\ell(c)$, an unsupported-fragment predicate $u_\ell(c)$, and an applicability predicate $\pi_\ell(c)$.  The predicates are recomputed from the policy and request material; they are not trusted labels from the generator.

\begin{equation}
\mathsf{status}(c)=
\begin{cases}
\Rej, & \rho_\ell(c)\\
\Unsup, & \neg\rho_\ell(c)\wedge u_\ell(c)\\
\App, & \neg\rho_\ell(c)\wedge\neg u_\ell(c)\wedge\pi_\ell(c)\\
\Rej, & \text{otherwise.}
\end{cases}
\end{equation}

For an applicable row, each declared adapter produces a before/follow-up decision pair $(d,d')$.  The row passes when the law invariant $I_\ell(d,d')$ holds.  It fails only when the row is applicable, supported, executable, and violates the invariant.  This is a narrow definition, but it is the right one for evaluator regression.  Unsupported rows should not be reported as passes; invalid transformations should not become failures; seeded drifts should not be mixed with primary evidence.

\begin{table}[t]
\centering
\caption{Terminal accounting in the current replay.}
\label{tab:accounting}
\begin{tabular}{p{0.37\linewidth}r p{0.31\linewidth}}
\toprule
Row family & Count & Meaning \\
\midrule
Applicable obligations & \LGPrimaryObligations{} & Counted backend-law checks \\
Primary passes & \LGPrimaryPasses{} & Law invariant held in counted rows \\
Primary failures & \LGPrimaryFailures{} & Real violations observed in counted rows \\
Rejected transformations & \LGRejected{} & Invalid law candidates, not counted \\
Unsupported fragments & \LGUnsupported{} & Outside declared adapter boundary \\
Seeded-drift rows & \LGSeededRows{} & Positive controls, kept separate \\
Minimized witnesses & \LGCounterexamples{} & Compact replay material for killed controls \\
\bottomrule
\end{tabular}
\end{table}

The conservative denominator view is included to make overfitting visible.  If rejected and unsupported candidates are counted against the run, the pass rate is \LGAllCandidatePassRate{}\%, not 100\%.  That number is not the primary metric because rejected and unsupported candidates are not valid evaluator obligations.  It is a pressure test: it tells the reader that the paper is not hiding the size of the excluded boundary.

\section{Law Kernel and Predicate Soundness}
The law catalog has \LGRelations{} relation families.  Each relation family contains a law statement, a side condition, a rejection rule, an unsupported-fragment rule, an invariant, and a minimization recipe.  The most important design choice is that the predicate must be semantic.  It may use syntax to identify the edited material, but admission to the denominator depends on recomputed reachability, matching, closure, dominance, or scope equivalence.

Deny-aware principal monotonicity illustrates the need for this discipline.  A naive law might say that substituting a principal with a role-closure superset preserves an allow decision.  Under deny-overrides semantics this is false: the larger closure can add a newly matching deny.  LatticeGuard therefore admits principal-order rows only when the follow-up closure is a superset and no newly matching deny interferes.  The side condition is part of the executable predicate, not a prose exception.

\begin{equation}
\begin{aligned}
\mathsf{PA}(P,q,P',q') \Rightarrow
& \mathsf{cl}_{P}(p) \subseteq \mathsf{cl}_{P'}(p') \\
& \wedge\ \neg \exists r.\ \mathsf{deny}(r)\wedge \mathsf{match}_{P'}(r,q')\\
& \wedge\ \Eval_{\Core}(P,q)=\Allow.
\end{aligned}
\end{equation}

When the side condition holds, the expected invariant is allow preservation.  When it does not hold, the row is rejected rather than counted.  The same pattern appears in other laws: a hierarchy closure row must preserve request-relevant reachability; an alpha-renaming row must not touch the requested atoms; a scope split/merge row must preserve the matched action/resource set; a rule-order row must avoid priority-sensitive fragments unless the law explicitly models priority.

\begin{table}[t]
\centering
\caption{Law-catalog evidence.}
\label{tab:lawkernel}
\begin{tabular}{p{0.36\linewidth}r p{0.31\linewidth}}
\toprule
Evidence & Count & Role \\
\midrule
Relation families & \LGRelations{} & Frozen law catalog \\
Theorem rows & 12 & Human-readable law statements \\
Proof obligations & 48 & Side-condition checks \\
Bounded core cases & \LGModelCheckCases{} & Finite falsification search \\
Bounded core failures & \LGModelCheckFailures{} & Unsound relations found \\
Independent kernel cases & \LGMechanizedKernelCases{} & Adapter-independent replay \\
Kernel failures & \LGMechanizedKernelFailures{} & Cross-check failures \\
\bottomrule
\end{tabular}
\end{table}

These checks are not presented as a complete theorem for all possible authorization languages.  They are a falsification-oriented guardrail for this study's law catalog.  If a future relation is added, it should pass the same style of side-condition checks before its generated candidates can affect any primary denominator.  That is the scientific point: the method makes law admission reviewable.

Table~\ref{tab:predicateobjects} spells out what a reviewer can inspect for each family.  The table is intentionally about predicates rather than implementation names.  A relation becomes scientific evidence only when the admitted row, the rejected row, and the minimized failure all have a stable semantic object.  This is also why the law catalog is not presented as a list of feature toggles.  The central claim is that relation families carry their own admission contracts.

\begin{table}[t]
\centering
\caption{Predicate objects used to keep relation families reviewable.}
\label{tab:predicateobjects}
\begin{tabular}{p{0.25\linewidth}p{0.25\linewidth}p{0.25\linewidth}}
\toprule
Relation family & Predicate object & Rejection boundary \\
\midrule
Default deny & Absence of reachable permit & Newly reachable allow \\
Deny dominance & Matching forbid over allow & Unmatched deny or allow \\
Forbid precedence & Dominating forbid slice & Backend lacks forbid support \\
Hierarchy closure & Request-relevant reachability & Closure edge changes query slice \\
Principal order & Closure superset plus no new deny & New deny becomes reachable \\
Resource order & Resource closure comparison & Resource path touches request \\
Alpha renaming & Injective off-slice mapping & Queried atom is renamed \\
Scope split/merge & Equivalent covered action set & Partition changes matched action \\
Shadowed removal & Dominated rule witness & Removed rule can decide request \\
Rule-order invariance & Priority-insensitive fragment & Priority-sensitive backend row \\
Monotone allow & Permit-preserving extension & Deny side condition violated \\
Unsupported fragment & Adapter support predicate & Outside declared fragment \\
\bottomrule
\end{tabular}
\end{table}

This table also clarifies why the bounded evidence is useful even though it is not a universal proof.  The finite search does not certify every future policy language extension.  It does something more local and reviewable: it checks that the predicates used in this run do not admit the known small countermodels for their own declared relation families.  If a relation fails this check, the right response is not to hide the row but to strengthen or remove the relation before it contributes to a denominator.

\section{Adapter and Corpus Boundaries}
The counted run uses two native evaluator roles and one normalized decision harness.  Casbin and Cedar are exercised through native evaluator semantics for their admitted fragments.  OPA/Rego is exercised as a normalized-decision fragment over a law object whose hierarchy closure is supplied by LatticeGuard's reference core.  That boundary matters.  The OPA rows show that the normalized allow/deny decision logic can be driven through Rego under a declared closure contract; they do not claim that recursive hierarchy closure is implemented natively by the Rego policy in this study.

\begin{table}[t]
\centering
\caption{Adapter roles in the evaluated fragment.}
\label{tab:adapterroles}
\begin{tabular}{p{0.22\linewidth}p{0.28\linewidth}p{0.30\linewidth}}
\toprule
Evaluator role & Counted semantics & Boundary \\
\midrule
Casbin & Native model, effect, and role matching & Priority-sensitive fragments need explicit support \\
Cedar & Native permit/forbid decision & Richer entity constraints need new contracts \\
OPA/Rego & Normalized allow/deny decision over supplied closure & Not a native recursive-closure claim \\
\bottomrule
\end{tabular}
\end{table}

The corpus is also stratified.  Documentation-derived slices and upstream examples provide public semantic anchors.  Native canonical slices preserve linkage to real backend fixtures through native self-tests.  Stress witnesses deliberately combine relation families that public examples rarely cover densely.  The generated law probe exercises canonical corner cases.  Collapsing these strata into one benchmark count would make the study look larger while saying less about external validity.

\begin{table}[t]
\centering
\caption{Subject strata and evidential roles.}
\label{tab:strata}
\begin{tabular}{p{0.34\linewidth}r p{0.34\linewidth}}
\toprule
Stratum & Count & Primary role \\
\midrule
Documentation-derived slices & \LGOfficialDocSources{} & Public semantics and examples \\
Upstream examples & \LGUpstreamExampleSources{} & Ecosystem-grounded policies \\
Native canonical slices & \LGNativeCanonicalSources{} & Backend fixture linkage \\
Semantic stress witnesses & \LGStressWitnessSources{} & Dense law-boundary coverage \\
Generated law probe & \LGGeneratedSources{} & Canonical corner case \\
\bottomrule
\end{tabular}
\end{table}

The native self-test layer contains \LGNativeSelftests{} raw decisions with \LGNativeFailures{} failures before normalization.  These self-tests do not prove that every native language feature is covered.  They do prevent a weaker failure mode: importing native material, simplifying it into a canonical form, and then forgetting whether the original evaluator would have made the same decisions.  In that sense, native self-tests are a provenance check for the corpus, not an additional headline result.

\begin{table}[t]
\centering
\caption{Boundary decisions for external validity.}
\label{tab:boundarydecisions}
\begin{tabular}{p{0.25\linewidth}p{0.25\linewidth}p{0.25\linewidth}}
\toprule
Reviewer question & Current decision & Why it is conservative \\
\midrule
Are all language features claimed? & No & Unsupported fragments remain outside counts \\
Are generated probes dominant? & No & Corpus strata are reported separately \\
Is OPA a full hierarchy engine? & No & Closure is supplied by the reference core \\
Are native examples decorative? & No & Self-tests bind native decisions first \\
Can public docs alone explain results? & No & Upstream-only control is measured separately \\
Can one adapter define truth? & No & Backend decisions are checked against reference semantics \\
\bottomrule
\end{tabular}
\end{table}

The boundary decisions are deliberately easy to audit.  A broader claim would be less useful here: if every unsupported feature were silently folded into the denominator, a clean run could be produced by excluding difficult semantics after the fact.  LatticeGuard instead makes exclusions visible.  The resulting study is narrower, but the counted obligations have a defensible interpretation.

\section{Denominator and Overfitting Checks}
A perfect primary pass rate is suspicious unless the denominator is inspectable.  The current replay therefore records \LGValidityBoundaryChecks{} validity-boundary checks across applicability accounting, unsupported fragments, OPA normalization, holdout coverage, seeded-drift sensitivity, corpus strata, source-boundary summaries, and deployed-tool head-to-head comparison.  These checks are designed to catch three families of overfitting: row selection after seeing outcomes, synthetic-source inflation, and denominator manipulation.  The \LGAllCandidatePassRate{}\% pressure view is intentionally pessimistic: it keeps \LGPrimaryPasses{} passed obligations in the numerator while counting rejected and unsupported candidates against the denominator rather than silently hiding them.

\begin{table}[t]
\centering
\caption{Validity-boundary checks.}
\label{tab:validityboundary}
\begin{tabular}{p{0.38\linewidth}p{0.48\linewidth}}
\toprule
Risk & Guardrail \\
\midrule
Invalid rows inflate pass rate & Rejection predicate recomputed before accounting \\
Unsupported rows disappear & Unsupported count remains visible and excluded \\
OPA boundary is overstated & \LGOPANormalizedRows{} normalized-decision rows reported under closure contract \\
Synthetic rows dominate story & Public/native/stress/generated strata remain separate \\
Failure path is inert & \LGSeededKillRate{}\% seeded positive-control kill rate \\
Baseline comparison is only conceptual & \LGDeployedHeadToHeadRows{} same-stream deployed-tool comparisons \\
Outcome tuning shapes evidence & \LGHoldoutObligations{} outcome-independent holdout obligations \\
Zero failures hide denominator pressure & \LGAllCandidatePassRate{}\% conservative pass-rate view \\
\bottomrule
\end{tabular}
\end{table}

The holdout evidence is intentionally outcome-independent.  It is not a new headline benchmark and does not let the paper claim broader language coverage.  Its role is to check that the admission and replay machinery can be applied to material that was not selected to explain the main primary pass count.  If the holdout failed admission checks or produced inconsistent accounting, the zero-primary-failure result would be much harder to trust.

The seeded-drift kill rate also needs careful interpretation.  The current value, \LGSeededKillRate{}\%, is not intended to be high for its own sake.  A positive-control family should exercise different relation failures without pretending that every seeded row is an equally realistic production bug.  A very high kill rate could itself become suspicious if it meant the mutants were trivial.  A zero kill rate would mean the failure path is inert.  The current result shows that the pipeline detects law-relevant semantic drifts while retaining a conservative interpretation.

\begin{table}[t]
\centering
\caption{Denominator pressure views.}
\label{tab:denominatorpressure}
\begin{tabular}{p{0.35\linewidth}r p{0.32\linewidth}}
\toprule
View & Value & Interpretation \\
\midrule
Primary applicable obligations & \LGPrimaryObligations{} & Valid evaluator-law rows \\
Rejected transformations & \LGRejected{} & Invalid candidates kept visible \\
Unsupported fragments & \LGUnsupported{} & Outside declared adapter support \\
All-candidate pass view & \LGAllCandidatePassRate{}\% & Pressure view, not headline metric \\
Holdout obligations & \LGHoldoutObligations{} & Outcome-independent replay surface \\
Validity-boundary checks & \LGValidityBoundaryChecks{} & Anti-overfitting gate count \\
Seeded-drift rows & \LGSeededRows{} & Separate positive-control stream \\
\bottomrule
\end{tabular}
\end{table}

The important feature of Table~\ref{tab:denominatorpressure} is that the paper does not ask the reader to accept a perfect pass rate in isolation.  The rejected, unsupported, holdout, and seeded streams remain numerically adjacent to the primary pass count.  That adjacency prevents a common overfitting pattern: selecting the denominator after seeing which candidates pass.  A future extension can improve coverage, but it must still preserve the same accounting separation.

\section{Counterexamples and Baselines}
When a seeded drift violates a law, LatticeGuard reduces the case using a relation-specific minimizer.  The minimized witness keeps the law name, the before/follow-up decision pair, the request material, and the predicate witness.  It removes irrelevant policy material only when the same predicate and invariant violation remain.  This is important because a large failing policy is easy to dismiss as an incidental artifact of generation.  A small law-named witness is closer to the object a maintainer needs after a release regression.

\begin{table}[t]
\centering
\caption{Why relation-specific minimization matters.}
\label{tab:minimization}
\begin{tabular}{p{0.31\linewidth}p{0.55\linewidth}}
\toprule
Failure family & Material that must survive minimization \\
\midrule
Deny dominance & Matching allow, matching deny, and request slice \\
Hierarchy closure & Edge witness and closure-relevant request \\
Principal monotonicity & Closure superset and no-new-deny side condition \\
Shadowed removal & Dominating deny and removed allow relation \\
Alpha renaming & Injective off-slice mapping and unchanged request atoms \\
Scope split/merge & Equivalent action/resource partition \\
\bottomrule
\end{tabular}
\end{table}

The baseline design is framed as competing explanations rather than as a generic feature checklist.  If native point tests and OPA-style decision tests were enough, before/follow-up law obligations would add little.  If cross-adapter disagreement were enough, law-specific applicability would be unnecessary.  If release-pair comparison alone were enough, future version pairs would not need predicate admission or replayable witnesses.  If random valid-looking mutations were enough, rejection accounting would not matter.  If a single relation family explained all seeded drifts, the catalog could be simplified.  The deployed crosswalk records \LGDeployedCrosswalkRows{} familiar idioms, \LGDeployedNativeSelftests{} native self-tests, and \LGDeployedCrossAdapterRows{} comparable differential rows before the component ablations are interpreted.  The deployed head-to-head then evaluates point-test and differential idioms as same-stream decision comparisons rather than as adapter-specific test-framework integration: seeded point checks report \LGSeededPointMismatches{} decision mismatches, \LGSeededHeadToHeadOverlap{} of those rows become law-level kills after predicate admission, and the point-only remainder stays visible outside the law-failure claim.  These alternatives are evaluated over the same frozen evidence stream, which keeps the comparison about oracle design rather than input selection.

\begin{table}[t]
\centering
\caption{Alternative explanations tested by the controls.}
\label{tab:baselines}
\begin{tabular}{p{0.34\linewidth}p{0.52\linewidth}}
\toprule
Alternative & What would have to be true \\
\midrule
Native and OPA point tests & Same-stream decision checks already explain law failures \\
Cross-adapter only & Disagreement alone explains semantic regressions \\
Release-pair precheck & Version comparison alone is enough for regression evidence \\
Random valid mutation & Syntactic validity is enough for metamorphic testing \\
No rejection discipline & Invalid transformations do not affect conclusions \\
No unsupported accounting & Skipped fragments cannot bias the result \\
Single relation catalog & One law family covers the semantic space \\
No minimization & Detection alone is sufficient for maintainers \\
No bounded checker & Empirical passes establish relation soundness \\
\bottomrule
\end{tabular}
\end{table}

This comparison does not claim that LatticeGuard dominates every possible authorization-testing tool on every task.  It claims something narrower and stronger: for regression oracles built from access-control laws, the critical difference is whether the expected relation is admitted by a recomputed predicate before evaluation and whether failures reduce to law-named witnesses.  Point checks can be stricter at the decision-table level and still be less explanatory at the law level; the baselines are chosen to attack exactly that claim.

\begin{table}[t]
\centering
\caption{Positive-control interpretation.}
\label{tab:positivecontrol}
\begin{tabular}{p{0.25\linewidth}p{0.25\linewidth}p{0.25\linewidth}}
\toprule
Seeded drift family & Expected witness & Reviewer use \\
\midrule
Allow stripping & Permit disappears from matched slice & Checks failure path is live \\
Deny stripping & Dominating deny disappears & Checks deny precedence logic \\
Hierarchy deletion & Closure edge is removed & Checks reachability predicates \\
Shadowing change & Dominated rule becomes visible & Checks minimizer keeps cause \\
Priority disturbance & Rule order changes decision & Checks unsupported boundaries \\
OPA closure mismatch & Supplied closure differs from policy object & Checks adapter contract \\
\bottomrule
\end{tabular}
\end{table}

The positive-control stream is not used to inflate the main pass rate.  It is a calibration device.  If the paper reported zero primary failures and no seeded controls, an obvious reviewer concern would be that the oracle is incapable of observing a violation.  If the seeded rows were mixed into the primary denominator, the opposite concern would arise: the paper would be counting artificial failures as deployment evidence.  Keeping the streams separate is what makes both readings possible without confusing them.

\begin{table}[t]
\centering
\caption{How the ablations attack the contribution.}
\label{tab:ablationattacks}
\begin{tabular}{p{0.27\linewidth}p{0.25\linewidth}p{0.23\linewidth}}
\toprule
Ablation & Attack being tested & Expected consequence \\
\midrule
No predicate admission & Law names are enough & Invalid rows enter evidence \\
No rejection ledger & Exclusions do not matter & Denominator becomes opaque \\
No unsupported accounting & Adapter gaps are harmless & Unsupported rows masquerade as passes \\
No native material & Synthetic examples suffice & External-validity anchor weakens \\
No bounded checker & Empirical rows prove laws & Unsound relation can enter run \\
No minimization & Detection alone is useful & Witnesses become hard to inspect \\
Single relation only & One law covers regressions & Multi-causal drifts are missed \\
\bottomrule
\end{tabular}
\end{table}

These ablations are stronger than a no-tool comparison because each one is allowed to explain away the same frozen evidence from a different angle.  The evaluation therefore asks whether the law-to-obligation design is necessary, not merely whether a named system can produce more rows than hand-written examples.

\section{Reproducibility and Residual Risks}
The replay surface is deterministic, CPU-local, offline, and anonymous.  The reported numeric claims are generated from the replay evidence, and the same evidence is checked for claim consistency, relation coverage, adapter/reference agreement, dependency closure, source linkage, and anonymization.  Operational details are kept separate from the paper prose so that the manuscript reads as a research contribution rather than as a laboratory notebook, but the release is still intended to let reviewers recompute the reported quantities.

\begin{table}[t]
\centering
\caption{Residual risks and current mitigations.}
\label{tab:residual}
\begin{tabular}{p{0.36\linewidth}p{0.50\linewidth}}
\toprule
Risk & Mitigation and boundary \\
\midrule
Law catalog is incomplete & New laws require predicates, side conditions, and replay evidence \\
OPA scope is overstated & Normalized-decision role and supplied closure are reported explicitly \\
Deployed baselines are argumentative & Same-stream head-to-head separates point mismatches from law failures \\
Stress witnesses dominate corpus & Strata are reported separately and not treated as prevalence evidence \\
Zero failures look overfit & Rejections, unsupported rows, holdout, seeded drifts, and denominator pressure are visible \\
Formal support is bounded & Bounded replay is presented as falsification support, not universal proof \\
External systems evolve & Future release-pair claims must re-enter through the same contracts \\
\bottomrule
\end{tabular}
\end{table}

The strongest remaining limitation is external validity.  Authorization languages contain features outside this study's core: obligations, environmental conditions, relationship tuples, recursive data structures, schema evolution, and service-level caching.  LatticeGuard's answer is not to claim those features for free.  A new feature must enter as a declared relation contract, adapter support rule, unsupported-fragment boundary, and replayable evidence stream.  That design may look slower than unrestricted generation, but it is what prevents the denominator from becoming a hidden modeling decision.

The second limitation is that primary failures are not observed in the current counted run.  The paper treats this as a regression-assurance result over declared fragments, not as a discovery claim.  In domains where current implementations already satisfy the evaluated laws, the value of a regression oracle is that it creates an inspectable baseline for future edits.  The seeded controls and minimizers show that the failure path is not decorative; a future real violation would enter the same schema and reduce to a replayable witness.

\begin{table}[t]
\centering
\caption{Claim-to-evidence reading path for a skeptical reviewer.}
\label{tab:readingpath}
\begin{tabular}{p{0.26\linewidth}p{0.25\linewidth}p{0.24\linewidth}}
\toprule
Concern & Evidence to inspect & What would fail \\
\midrule
Perfect pass rate is overfit & Rejected, unsupported, holdout, seeded streams & Hidden denominator manipulation \\
OPA scope is overstated & Normalized-decision role and closure contract & Native-recursion overclaim \\
Baselines are ornamental & Component-removal controls & Feature checklist comparison \\
Formalism is decorative & Bounded and kernel law checks & Unsound relation admission \\
Corpus is synthetic & Public, upstream, native, stress strata & Single-source benchmark story \\
Failures are not actionable & Relation-specific minimization & Large unreplayable witnesses \\
Artifact is stale & Claim traceability and replay gates & Paper-number drift \\
\bottomrule
\end{tabular}
\end{table}

The reading path is meant to make rejection easier if the evidence does not support the claim.  Each row names a way the paper could be wrong and the evidence surface that would reveal it.  This is the standard the project should preserve after publication as well: adding a new adapter, language feature, or relation family should add a new auditable boundary, not only a larger count.

The review package is therefore best read as a regression-oracle construction rather than a benchmark leaderboard.  Its contribution is not that every authorization language has been covered.  Its contribution is that an evaluator-law claim can be admitted, rejected, minimized, bounded, and replayed without letting invalid transformations become silent evidence.  That is the property future release-pair studies can reuse.

\begin{table}[t]
\centering
\caption{Acceptance criteria for future extensions.}
\label{tab:extensioncriteria}
\begin{tabular}{p{0.32\linewidth}p{0.48\linewidth}}
\toprule
Extension type & Required evidence before counting \\
\midrule
New relation family & Law statement, side condition, rejection predicate, invariant, and minimizer \\
New backend fragment & Adapter support rule, reference agreement, and unsupported-fragment boundary \\
New corpus source & Public or native provenance stratum and self-test or semantic anchor \\
New release-pair study & Frozen before/after material and replayable claim ledger \\
New seeded drift & Declared failure family and relation-specific minimized witness \\
New baseline & Competing explanation stated before looking at outcomes \\
New formal check & Bound, domain, and relation family stated explicitly \\
\bottomrule
\end{tabular}
\end{table}

These criteria are stricter than ordinary benchmark growth.  A larger corpus or another backend is valuable only if it preserves the claim boundary.  The extension must say what is newly supported, what remains unsupported, and how a violation would reduce to a concrete witness.  Otherwise growth would make the artifact look stronger while making the scientific interpretation weaker.

This is also the reason the release keeps one current evidence surface rather than multiple named revisions.  Reviewer-facing claims should refer to the same frozen replay and the same paper-visible numbers.  If a later run changes any denominator, support count, or failure family, the correct response is to replace the evidence surface and update the claims, not to leave parallel versions for reviewers to reconcile.

\end{document}
